\section{Probabilistic Model of Building Fire Risk}
\label{sec:model}

We treat the fire risk of a residential building $b$ not as a categorical score
but as a \emph{posterior distribution} over the coupled processes of ignition,
fire development, evacuation failure, and consequence severity, updated as
evidence arrives. This section formalises that view.

\subsection{Risk as expected loss}

Let $I_b(t)$ denote the event that a fire ignites in building $b$ during
year $t$, and let $C_b(t)$ denote the resulting consequence (loss). We define the
annualised risk as the expected loss,
\begin{equation}
  R_b(t) \;=\; P\!\big(I_b(t)\big)\,\cdot\,
  \mathbb{E}\!\big[\,C_b(t)\,\big|\,I_b(t)\,\big],
  \label{eq:risk}
\end{equation}
i.e.\ the product of ignition probability and expected consequence given
ignition. This mirrors the decomposition used in probabilistic safety analysis
and in Bayesian-network dwelling-fire models
\citep{matellini2018threepart,estradalugo2019grenfell}.

\subsection{Consequence decomposition}

Consequence is not a single quantity: a fire can destroy property without harming
anyone, and can injure without spreading beyond the room of origin. We therefore
decompose expected consequence as a \emph{sum over loss channels},
\begin{equation}
  \mathbb{E}\!\big[C_b \mid I_b\big] \;=\;
  \mathbb{E}\!\big[C^{\mathrm{prop}}_b \mid I_b\big] +
  \mathbb{E}\!\big[C^{\mathrm{inj}}_b \mid I_b\big] +
  \mathbb{E}\!\big[C^{\mathrm{fat}}_b \mid I_b\big],
  \label{eq:consequence-channels}
\end{equation}
and, within each channel, marginalise over the fire-development pathway rather
than forcing every loss through a single chain. Writing $S$ for the extent of
spread (an ordered variable: contained to item, room of origin, beyond room,
beyond floor) and $E_f$ for evacuation failure,
\begin{equation}
  P\!\big(\text{casualty} \mid I_b\big) \;=\;
  \sum_{s}\;\sum_{e \in \{E_f, \bar{E}_f\}}
  P\!\big(\text{casualty} \mid e, s, I_b\big)\,
  P\!\big(e \mid s, I_b\big)\,
  P\!\big(s \mid I_b\big),
  \label{eq:consequence}
\end{equation}
which preserves the causal ordering ignition $\to$ spread $\to$ egress $\to$ harm
while permitting casualties without major spread (e.g.\ smoke inhalation or
fire-fighting attempts in the room of origin, the $s{=}\text{room}$,
$e{=}\bar{E}_f$ terms) and property loss without evacuation failure. The dominant
term for fatalities is the one highlighted by the multiplicative chain
$P(S{\geq}\text{room}\mid I)\,P(E_f\mid S,I)\,\mathbb{E}[L\mid E_f,S,I]$, which we
retain as the leading-order approximation used in the worked examples; Equation~%
\eqref{eq:consequence} is the full object. Each factor depends on the building's
physical and occupancy characteristics; the egress factors are where the
physics-based consequence surrogate of Section~\ref{sec:surrogate} enters,
and the property channel is valued directly from incident-damage statistics.

\paragraph{Why this factorisation.} Four modelling choices in
Equation~\eqref{eq:consequence} warrant defence. First, spread is marginalised
\emph{before} egress because the two orderings coincide: compartmentation failure
is physically upstream of tenability loss, so conditioning harm on spread respects
the causal chain, and spread is also the margin the national incident record
actually reports, so the factorisation matches an observable quantity rather than
a latent one. Second, evacuation is kept binary ($E_f$ versus $\bar E_f$) because
the record carries no evacuation \emph{timing}, only its outcome; a finer egress
state would be prior, not data. Third, suppression is not a separate factor:
sprinkler coverage in the existing English dwelling stock is of the order of one
per cent, so it is negligible as a population term, and brigade intervention acts
\emph{after} the spread outcome that defines our margin, so its effect is already
absorbed into the empirically estimated $P(s\mid I_b)$ rather than double-counted.
Fourth, we do not carry a fully joint latent $(\text{spread},\text{occupancy},
\text{detection})$ object here: that model is written down in
Section~\ref{sec:model-inspection}, but it is not identified from a single
cross-section, so the factorised form of Equation~\eqref{eq:consequence} is the
part the present data support.

\subsection{Latent building state and evidence}

The factors in Equations~\eqref{eq:risk}--\eqref{eq:consequence} are governed by
a \emph{latent building state}
\begin{equation}
  \begin{aligned}
  z_b \;=\; \{\, &\text{compartmentation},\ \text{fuel load},\
  \text{alarm reliability},\ \text{suppression},\\
  &\text{occupancy vulnerability},\ \text{maintenance quality},\
  \text{electrical risk} \,\},
  \end{aligned}
  \label{eq:latent}
\end{equation}
which is never observed directly. Instead we observe \emph{evidence}
\begin{equation}
  \begin{aligned}
  x_b \;=\; \{\, &\text{inspection findings},\ \text{public building attributes},\\
  &\text{incident history},\ \text{occupancy proxies},\
  \text{local covariates} \,\}.
  \end{aligned}
  \label{eq:evidence}
\end{equation}
Inference proceeds in two steps: recover the posterior over latent state,
\begin{equation}
  P\!\big(z_b \mid x_b\big) \;\propto\;
  P\!\big(x_b \mid z_b\big)\,P\!\big(z_b\big),
  \label{eq:posterior-z}
\end{equation}
and propagate it through the risk model to obtain the posterior risk
distribution,
\begin{equation}
  P\!\big(R_b \mid x_b\big) \;=\;
  \int P\!\big(R_b \mid z_b\big)\,P\!\big(z_b \mid x_b\big)\,\mathrm{d}z_b .
  \label{eq:posterior-R}
\end{equation}
Equation~\eqref{eq:posterior-R} is the central object of the framework: \emph{risk
is a distribution, not a score}. Its spread quantifies uncertainty arising from
missing or unreliable evidence, and its credible intervals and missing-evidence
flags are the operational outputs consumed by FRA workflows
(Section~\ref{sec:results}).

\paragraph{Implemented vs proposed.} This paper implements calibrated components
of Equations~\eqref{eq:risk}--\eqref{eq:consequence} --- a small-area occurrence
model for $P(I)$, record-level consequence models, and the consequence simulator ---
composed with propagated uncertainty in Sections~\ref{sec:results-synthesis}
and~\ref{sec:results-decision}. The joint latent-state posterior of
Equations~\eqref{eq:posterior-z}--\eqref{eq:posterior-R}, with the variables of
Equation~\eqref{eq:latent}, is specified in full but only partially instantiated:
Section~\ref{sec:results-latent} fits its measurement model on 1{,}441 published
fire risk assessments, while the transition dynamics remain unidentified until
repeated assessment waves exist.

\paragraph{Predictive association vs intervention effect.} Throughout we
distinguish the \emph{predictive} quantity $P(Y{=}1 \mid x, I)$, which record-level
regressions estimate, from the \emph{interventional} contrast
$P\!\big(Y^{\mathrm{do(alarm=working)}}{=}1\big) -
 P\!\big(Y^{\mathrm{do(alarm=absent)}}{=}1\big)$. Observational fire data conflate
the two: alarm operation, occupancy, and casualty are mutually selected, since an
alarm operates, and a casualty can occur, only when someone is present. Predictive
associations are therefore used for risk ranking, while intervention effects are
taken from the mechanism-informed simulator, whose alarm effect is interventional
by construction (Section~\ref{sec:surrogate}). The distinction is drawn here,
before any results, because Section~\ref{sec:results} reports estimates whose
observational sign differs from the interventional one.

\subsection{Inspection as noisy observation}
\label{sec:model-inspection}

We treat an inspection not as ground truth but as a noisy, intermittent
measurement of the latent state, in the spirit of detector/measurement-error
modelling. For a latent attribute $z_b^{(k)}$ (e.g.\ compartmentation integrity),
an inspection at time $t$ yields an observation $o_b^{(k)}(t)$ through an
observation model
\begin{equation}
  o_b^{(k)}(t) \;\sim\; P\!\big(o \mid z_b^{(k)},\, \phi_k \big),
  \label{eq:obs-model}
\end{equation}
where $\phi_k$ captures inspection sensitivity, specificity, and inspector
variability. Absence of a recent inspection simply means no update to
$P(z_b^{(k)}\mid x_b)$, so uncertainty grows with time since last observation.

\paragraph{Latent-state dynamics, and what the data instantiate.} Indexing the
state by time, $z_{b,t}$ follows a first-order state space with interventions,
\begin{align}
  z_{b,t} &\;\sim\; p\big(z_{b,t} \mid z_{b,t-1},\, u_{b,t}\big)
  &&\text{(slow drift; jumps at refurbishment/intervention $u_{b,t}$)},
  \label{eq:ss-transition}\\
  o_{b,t} &\;\sim\; p\big(o_{b,t} \mid z_{b,t},\, \phi\big)
  &&\text{(inspection observations, Eq.~\eqref{eq:obs-model})},
  \label{eq:ss-obs}\\
  y_{b,t} &\;\sim\; p\big(y_{b,t} \mid z_{b,t},\, x_{b,t}\big)
  &&\text{(incident outcomes: the models of Section~\ref{sec:results})},
  \label{eq:ss-outcome}
\end{align}
so a new inspection updates $P(z_{b,t}\mid o_{b,1:t}, y_{b,1:t}, x_{b,1:t})$ by
standard filtering, and drift between inspections widens it. The assessor
parameters $\phi$ --- detection sensitivity, specificity, between-assessor
variance --- are the empirically hardest piece. Section~\ref{sec:results-latent}
fits the measurement model of Equation~\eqref{eq:ss-obs} on 1{,}441 published fire
risk assessments across 653 Camden estates and \emph{bounds} detection sensitivity
from their recurrence structure; point identification of $\phi$ and of the
transition dynamics in Equation~\eqref{eq:ss-transition} needs repeated assessment
waves of the same buildings. We specify the model now because the repeated
assessments mandated under the Building Safety Act will identify these parameters;
the contribution here is the framework and the first empirical persistence
estimates, not a fitted latent state.

This framing also yields the \emph{expected value of inspection} (EVI): with actions
$a$ and losses $L$, acquiring observation $o_b$ before acting is worth
\begin{equation}
  \mathrm{EVI}(b) \;=\;
  \min_a \mathbb{E}\!\big[L(a) \mid x_b\big]
  \;-\;
  \mathbb{E}_{o_b}\!\Big[\min_a \mathbb{E}\!\big[L(a) \mid x_b, o_b\big]\Big]
  \;\;\geq\; 0,
  \label{eq:evi}
\end{equation}
the classic value-of-information quantity: it is largest where the posterior is
both consequential and undecided, and zero where the optimal action is already
determined. Section~\ref{sec:results-decision} instantiates
Equation~\eqref{eq:evi} for alarm-ownership inspection and shows the resulting
ranking differs materially from expected-risk ranking.

\subsection{Decision layer: allocating a prevention budget}
\label{sec:model-decision}

The posterior risk distributions exist to drive decisions, so the framework
includes the allocation problem explicitly. Let area $i$ contain $H_i$ dwellings,
let $n_i$ be the number visited under a prevention programme with budget $B$, and
let $b_i$ be the expected present-value averted loss per visited dwelling ---
composed from the posterior fire rate, the probability the dwelling lacks a
working alarm, and the interventional consequence reduction of
Section~\ref{sec:surrogate}. The allocation problem is
\begin{equation}
  \max_{n}\; \sum_i n_i\, b_i
  \qquad \text{s.t.} \quad \sum_i n_i \le B,
  \quad 0 \le n_i \le H_i .
  \label{eq:knapsack}
\end{equation}
Because every dwelling within an area shares the same marginal value $b_i \ge 0$,
Equation~\eqref{eq:knapsack} is a continuous knapsack with box constraints, and
the greedy rule the implementation uses --- fill areas in decreasing order of
$b_i$ until the budget is exhausted --- is exactly optimal.

\begin{proposition}[Greedy optimality of the allocation]
\label{prop:greedy}
Let $b_i \ge 0$ and index the areas so that $b_{(1)} \ge b_{(2)} \ge \cdots$.
In the continuous relaxation of Equation~\eqref{eq:knapsack} ($n_i \in [0,H_i]$,
$\sum_i n_i \le B$), the greedy allocation that fills areas in this order ---
$n_{(1)} = H_{(1)}, n_{(2)} = H_{(2)}, \dots$ until the running total reaches $B$,
filling the final boundary area only partially and setting the remainder to
zero --- attains the maximum of $\sum_i n_i b_i$.
\end{proposition}

\begin{proof}
A maximiser exists: the objective is linear and the feasible set compact. Let
$n^\ast$ be any maximiser and suppose it is not of the greedy form. Then there is
a pair $(j,k)$ with $b_j > b_k$, $n^\ast_j < H_j$ and $n^\ast_k > 0$: a
higher-value area not yet full while a lower-value one still draws budget. Moving
$\delta = \min(H_j - n^\ast_j,\; n^\ast_k) > 0$ from $k$ to $j$ leaves
$\sum_i n_i$ unchanged and both box constraints satisfied, so the result is
feasible, and it raises the objective by $\delta\,(b_j - b_k) > 0$ ---
contradicting the optimality of $n^\ast$. Every maximiser is therefore of the
greedy form, and all greedy-form allocations share the same objective value
(ties $b_j = b_k$ only permute budget between equal-value areas), so the greedy
allocation attains the maximum. (If $B \ge \sum_i H_i$ the budget binds nowhere
and $n_i = H_i$ is optimal because $b_i \ge 0$.)
\end{proof}

\noindent\emph{Remarks.} Were the per-area benefit strictly concave, $g_i(n_i)$
with $g_i'$ non-increasing, the KKT conditions of the same concave programme would
replace ``fill in order'' by \emph{water-filling}: interior areas equalise their
marginal benefit $g_i'(n_i) = \mu$ at the budget multiplier $\mu$, and the linear
case is the degenerate limit in which a constant marginal $g_i' = b_i$ forces each
area to fill completely before the next opens. Because visits are integer
dwellings, the deployed allocation rounds only the single boundary area, so the
gap between the integer optimum and the continuous one is at most the benefit of
one dwelling there, $\le \max_i b_i$ --- negligible against a budget of order
$10^5$ visits.

Uncertainty propagates by solving Equation~\eqref{eq:knapsack} per posterior draw,
giving credible intervals on the value of any policy; and replacing $b_i$ with
$\mathrm{EVI}_i$ of Equation~\eqref{eq:evi} yields the inspection-prioritisation
variant.

\paragraph{Propagating posterior uncertainty to the decision.} The pipeline is a
posterior $\pi(\theta \mid y)$, a loss $L(a,\theta)$ evaluated through the
simulator, and a plug-in decision that minimises the posterior expected loss. Two
guarantees govern how uncertainty in $\theta$ reaches the chosen action.

\begin{assumption}[Lipschitz loss]
\label{ass:lipschitz}
The action set $A$ is finite, and for each $a \in A$ the loss $L(a,\cdot)$ is
$L_a$-Lipschitz in $\theta$ on a set containing both the support of
$\pi(\theta \mid y)$ and the point estimate $\hat\theta$, i.e.\
$|L(a,\theta) - L(a,\theta')| \le L_a \|\theta - \theta'\|$. Write
$L_{\max} = \max_{a \in A} L_a$.
\end{assumption}

\begin{proposition}[Propagation and decision robustness]
\label{prop:propagation}
Write $\bar L(a) = \mathbb{E}_\pi[L(a,\theta) \mid y]$ for the posterior expected
loss and $a^\star = \arg\min_{a} \bar L(a)$ for the Bayes action.
\begin{enumerate}
\item[(i)] \textup{(Monte Carlo propagation.)} For $S$ posterior draws
$\theta_s \sim \pi(\cdot \mid y)$, the estimator
$\hat L_S(a) = S^{-1}\sum_{s=1}^{S} L(a,\theta_s)$ is unbiased for $\bar L(a)$ with
standard error $\sqrt{\mathrm{Var}_\pi[L(a,\theta)]\,/\,S}$. If the draws are
importance-weighted with self-normalised weights $w_s$ ($\sum_s w_s = 1$), the
estimator $\sum_s w_s\, L(a,\theta_s)$ has, by the delta method, variance
approximately $\mathrm{Var}_\pi[L(a,\theta)]\,/\,S_{\mathrm{eff}}$ with effective
sample size $S_{\mathrm{eff}} = 1/\sum_s w_s^2$ --- the importance-sampling
effective sample size of Kong (1992).
\item[(ii)] \textup{(Decision robustness.)} Let $\hat\theta$ be any point estimate
and $\hat a = \arg\min_{a} L(a,\hat\theta)$ the plug-in action. Under
Assumption~\ref{ass:lipschitz} the Bayes regret of acting on $\hat\theta$ obeys
\begin{equation}
  \bar L(\hat a) - \bar L(a^\star)
  \;\le\; 2\,L_{\max}\,\mathbb{E}_\pi\|\theta - \hat\theta\| .
  \label{eq:regret-bound}
\end{equation}
\end{enumerate}
\end{proposition}

\begin{proof}
Part (i) is the standard Monte Carlo variance, with the self-normalised
importance-sampling correction $S_{\mathrm{eff}} = (\sum_s w_s)^2 / \sum_s w_s^2 =
1/\sum_s w_s^2$ for normalised weights. For part (ii), write
$\varepsilon = \mathbb{E}_\pi\|\theta - \hat\theta\|$. For every $a$,
\begin{equation*}
  |\bar L(a) - L(a,\hat\theta)|
  = \big|\,\mathbb{E}_\pi[L(a,\theta) - L(a,\hat\theta)]\,\big|
  \le \mathbb{E}_\pi\,|L(a,\theta) - L(a,\hat\theta)|
  \le L_a\,\varepsilon \le L_{\max}\,\varepsilon .
\end{equation*}
Decompose the regret by inserting $L(\cdot,\hat\theta)$:
\begin{equation*}
  \bar L(\hat a) - \bar L(a^\star)
  = \underbrace{[\,\bar L(\hat a) - L(\hat a,\hat\theta)\,]}_{\le\, L_{\max}\varepsilon}
  + \underbrace{[\,L(\hat a,\hat\theta) - L(a^\star,\hat\theta)\,]}_{\le\, 0}
  + \underbrace{[\,L(a^\star,\hat\theta) - \bar L(a^\star)\,]}_{\le\, L_{\max}\varepsilon}.
\end{equation*}
The middle bracket is non-positive because $\hat a$ minimises $L(\cdot,\hat\theta)$;
the outer two are each at most $L_{\max}\varepsilon$ by the preceding display.
Summing gives Equation~\eqref{eq:regret-bound}.
\end{proof}

Part (ii) makes the value of sharp inference concrete: posterior contraction from
full-data inference (Section~\ref{sec:results}) shrinks
$\mathbb{E}_\pi\|\theta - \hat\theta\|$ and so tightens the regret bound directly,
meaning a point-estimate decision forfeits least precisely where the posterior is
already concentrated. Section~\ref{sec:results-decision} evaluates both allocation
variants against random and deprivation-based allocation.

\subsection{Rare events and calibration}
\label{sec:model-rare}

Building-level fire is a rare event: annual ignition probabilities are small and
severe outcomes rarer still \citep{mhclg2026firestats}. This has two
consequences for the model. First, estimation must borrow strength across
buildings via hierarchical/partial-pooling priors on $P(I_b)$ and on the
consequence factors, rather than fitting each building independently. Second,
\emph{calibration} is a first-class objective: predicted probabilities and
credible intervals must match observed frequencies at the population level.
We therefore evaluate the model with proper scoring rules and calibration
diagnostics (reliability curves, coverage of credible intervals) against
historical English incident statistics, treating the aggregate counts as the
frequency ground truth for the rare-event process (Section~\ref{sec:results}).

\subsection{From specification to fit: estimation and evaluation machinery}
\label{sec:model-estimation}

The results of Section~\ref{sec:results} instantiate the framework with three
fitted likelihoods; we state them fully here so that every fit is reproducible
from the paper alone.

\paragraph{Occurrence.} Dwelling-fire counts $y_{it}$ in LSOA $i$ and fiscal
year $t$ follow a hierarchical negative binomial in its mean--dispersion
(NB2) form,
\begin{align}
  y_{it} &\;\sim\; \mathrm{NB2}\big(\mu_{it},\, \phi\big),
  \qquad
  \mathrm{Var}[y_{it}] = \mu_{it} + \mu_{it}^2/\phi,
  \label{eq:nb2}\\
  \log \mu_{it} &\;=\; \log H_i \;+\; \alpha
  \;+\; \mathbf{x}_{it}^{\top}\boldsymbol\beta
  \;+\; b\,(t - \bar t)
  \;+\; u_{f(i)} \;+\; s_i ,
  \label{eq:occ-linpred}
\end{align}
with households $H_i$ as exposure offset, ten standardised covariates
$\mathbf{x}_{it}$, a linear year trend, service random intercepts
$u_f \sim \mathcal{N}(0, \sigma_u^2)$, and --- in the preferred M5
specification --- a BYM2 spatial effect
$s_i = \sigma_s\big(\sqrt{1-\rho}\; v_i + \sqrt{\rho}\; w_i/\sqrt{\kappa}\big)$
combining unstructured noise $v_i \sim \mathcal{N}(0,1)$ with an intrinsic
conditional autoregression $\mathbf{w} \sim \mathrm{ICAR}(\mathbf{A})$ on the
queen-contiguity LSOA adjacency $\mathbf{A}$, scaled by the graph's marginal
variance $\kappa$ so that $\sigma_s$ and the mixing weight $\rho$ are
interpretable. Priors are weakly informative:
$\alpha \sim \mathcal{N}(-6.5, 1.5^2)$ centred on the historical national
log-rate, $\beta_k \sim \mathcal{N}(0,1)$,
$b \sim \mathcal{N}(0, 0.5^2)$, $\sigma_u \sim \mathrm{HalfNormal}(0.5)$,
and $1/\phi \sim \mathrm{HalfNormal}(1)$, the last parameterised through the
inverse concentration so the Poisson limit sits at the prior mode's edge rather
than at a boundary singularity --- the reparameterisation that restores
tractable posterior geometry.

\paragraph{Consequence.} Record-level outcomes (serious spread; any casualty)
follow hierarchical logistic regressions,
\begin{equation}
  y_{jf} \sim \mathrm{Bernoulli}(p_{jf}), \qquad
  \mathrm{logit}\, p_{jf} = \alpha + \mathbf{x}_j^{\top}\boldsymbol\beta
  + u_f, \qquad u_f \sim \mathcal{N}(0, \sigma_u^2),
  \label{eq:cons-logit}
\end{equation}
with reference-coded factors for dwelling type, cause, ignition power,
occupancy, time of day, the four-level alarm factor, a standardised year trend
and the service-year response-time control; priors
$\alpha \sim \mathcal{N}(0, 5^2)$, $\beta \sim \mathcal{N}(0, 2.5^2)$,
$\sigma_u \sim \mathrm{HalfNormal}(1)$.

\paragraph{How the fits are performed.} Exact posteriors are drawn with the
No-U-Turn Sampler (NUTS), the adaptive Hamiltonian Monte Carlo variant that
simulates Hamiltonian trajectories on the log-posterior and eliminates
hand-tuned step counts. Hierarchical models of this type develop funnel-shaped
geometry at full data scale on which the NUTS step size collapses; we therefore
follow a two-track protocol. Track one: exact NUTS on stratified subsamples
($\sim$10{,}000 LSOAs spanning every service; $\sim$30{,}000 fires stratified
by fiscal year), 2--4 chains, 300--1{,}000 warmup and post-warmup draws,
accepted only with $\hat{R} \le 1.02$, bulk effective sample sizes above
$\sim$150 for every reported quantity, and zero divergent transitions. Track
two: the identical model on the \emph{full} data by automatic-differentiation
variational inference (ADVI), fitting a full-rank Gaussian
$q_{\lambda}(\boldsymbol\theta)$ by maximising the evidence lower bound
\begin{equation}
  \mathrm{ELBO}(\lambda) \;=\;
  \mathbb{E}_{q_\lambda}\big[\log p(\mathbf{y}, \boldsymbol\theta)\big]
  \;-\; \mathbb{E}_{q_\lambda}\big[\log q_\lambda(\boldsymbol\theta)\big],
  \label{eq:elbo}
\end{equation}
which serves as a point-estimate cross-check on all the data: coefficients must
agree with the subsample NUTS posterior within its credible intervals. Known
variational failure modes --- ridge-splitting between an intercept and its
random-effect mean, understated group-level predictive uncertainty --- are
checked for and reported wherever they occur.

Where exact NUTS is feasible on the full data, it supersedes both tracks: the
consequence models (logistic likelihood, no dispersion parameter) converge
cleanly on all 451{,}172 records and those posteriors are the ones reported.
The occurrence model does not admit full-data NUTS: on all 202{,}530 training
rows the negative-binomial posterior concentrates onto a ridge between the
intercept, the service-effect mean and the dispersion parameter that defeats
every standard reparameterisation --- non-centred random effects mix at
effective sample sizes of single digits, the centred form collapses the
adapted step size to $4\times10^{-6}$, and a QR-orthogonalised design (which
removes the deprivation-domain collinearity) still collapses to
$5\times10^{-7}$. Sampling difficulty is however a \emph{traversal} problem,
and importance sampling requires only \emph{coverage}: with 58 unconstrained
parameters and $2\times10^{5}$ observations the posterior is close to Gaussian
(Bernstein--von Mises), so the exact full-data posterior is recovered by
Pareto-smoothed importance sampling (PSIS) \citep{vehtari2024psis} over a
Laplace proposal. Concretely: the model is written in its \emph{centred}
parameterisation --- hopeless for NUTS, but the right coordinates here, since
each service's thousands of observations pin its effect directly and the
residual intercept--service-mean ridge is linear, which a Gaussian proposal
absorbs exactly through the Hessian; the mode and dense Hessian are computed
in unconstrained space; a multivariate Student-$t$ ($\nu=8$, scale 1.1) with
that mean and covariance proposes 50{,}000 draws; and the importance weights
$w_s \propto p(\boldsymbol\theta_s \mid \mathbf{y}) / q(\boldsymbol\theta_s)$
are stabilised by generalised-Pareto smoothing of the largest weights. The
smoothing diagnostic certifies the result: $\hat{k} = 0.33$, far below the 0.7
reliability threshold, with an importance-sampling effective sample size of
16{,}892 --- and the recovered posterior matches the subsample NUTS, ADVI and
maximum-likelihood (MLE) fits on every parameter. The two-track protocol is thus retained as
scaffolding and cross-check, while every headline posterior in this paper is
an exact full-data object.

\paragraph{How calibration is measured.} For discrete counts, central-interval
coverage over-covers structurally, so the primary diagnostic is the randomised
probability integral transform: for held-out count $y$ with predictive
cumulative distribution $F$,
\begin{equation}
  u \;=\; F(y-1) \;+\; v\,\big[F(y) - F(y-1)\big],
  \qquad v \sim \mathrm{Uniform}(0,1),
  \label{eq:pit}
\end{equation}
which is exactly $\mathrm{Uniform}(0,1)$ if and only if the predictive
distribution is calibrated; we report the Kolmogorov--Smirnov distance of $u$
from uniformity and the randomised coverage of central 50\%/90\% intervals.
Binary outcomes are scored with the Brier score
$\tfrac1n\sum_j (p_j - y_j)^2$, the mean log score, and the expected
calibration error $\mathrm{ECE} = \sum_m \tfrac{n_m}{n}\,
\lvert \bar p_m - \bar y_m \rvert$ over $M{=}10$ equal-width probability
bins, alongside coverage of service-level posterior-predictive count
intervals. All evaluation is out-of-time: the final one or two fiscal years
are never seen in fitting.

\paragraph{Fitting the measurement model.} The latent-state instantiation of
Section~\ref{sec:results-latent} operationalises
Equations~\eqref{eq:ss-obs}--\eqref{eq:ss-outcome} cross-sectionally: a scalar
latent deficit $\theta_e \sim \mathcal{N}(0,1)$ per estate $e$ loads on the
assessment indicators --- ordinal overall rating via a cumulative-logit link,
priority-A and priority-B action counts via Poisson links
$\log \mathbb{E}[\text{count}] = a_k + \lambda_k \theta_e$, and
deficiency-theme flags via Bernoulli-logit links --- while the structural
outcome is the estate's long-run fire count,
\begin{equation}
  y_e \sim \mathrm{NB2}\big(\mu_e, \phi\big), \qquad
  \log \mu_e = \log E_e + \alpha + \beta_\theta\, \theta_e
  + \beta_{\mathrm{IMD}}\, \mathrm{IMD}_e,
  \label{eq:measurement}
\end{equation}
with exposure $E_e$ (flats $\times$ years) and orientation pinned by
constraining the priority-A loading positive. The posterior update shown in
Figure~\ref{fig:latent} is then literal conditioning: the block's
$P(\theta_e \mid \text{FRA indicators})$ propagated through
Equation~\eqref{eq:measurement} to a fire-rate posterior, before and after the
indicators are observed.
